
\section{Selected Data Sources} \label{data}

Following the project guidelines to integrate at least two different data sources, this project integrates five distinct datasets containing spatial and administrative information. As detailed in the guidelines, these data sources were introduced within the context of the CRIMA project. Some information about them was documented in the Teams channel of the Data Profiling module, specifically located in the folder \inlinecode{data/crima} project. Specifically, the integration will involve the following provided data sources:
\begin{itemize}
  \item \textbf{\inlinecode{bz_infrastructure_line}:} This dataset provides geometric and categorical information regarding critical infrastructure lines. It contains multi-line geometries representing physical lines alongside classifications for infrastructure types, such as power lines and water pipes. The line dataset contains 47,551 rows detailing 11 types of line infrastructure.
  \item \textbf{\inlinecode{bz_infrastructure_node}:} This dataset contains point-like critical infrastructure information. It specifies the point geometries and infrastructure types for various nodes, such as water catchments, electrical substations, and communication infrastructure. The node dataset contains 4,223 rows representing 10 types of facilities.
  \item \textbf{\inlinecode{bz_hazard_landslide}:} This source supplies information regarding known landslide risk areas. It includes polygon geometries defining the spatial extent of the hazards, alongside categorical classifications of the landslide process types and level of danger. The landslide dataset contains 23,737 records categorized into four danger levels and five specific landslide process types.
  \item \textbf{\inlinecode{bz_hazard_avalanche}:} Similar to the landslide dataset, this source provides spatial boundaries and risk evaluations for avalanche-prone areas. It includes polygon geometries, avalanche process categorizations, and danger levels. The avalanche dataset includes 13,071 records, classified like the landslide dataset by danger severity and avalanche process type.
  \item \textbf{\inlinecode{bz_admin_municipality}:} This dataset provides the administrative context for the region. It contains 116 records including polygon geometries representing municipal boundaries, names in multiple languages, and district associations. This dataset allows the system to aggregate and assign hazard risks to specific municipalities.
\end{itemize}

Originally these datasets were derived from the Bolzano Province GeoKatalog. They provide the necessary geospatial geometries and categorical data required for the domain analysis. All the datasets were provided for this project in the form of GZip compressed CSV files. All datasets include geographic information stored in the Well-Known Text format. The files contain multiple data quality issues that will be identified and resolved as part of the data profiling part of this project.

